\documentclass[11pt]{article}
\usepackage[a4paper,margin=27mm]{geometry}
\usepackage[T1]{fontenc}
\usepackage{lmodern}
\usepackage{amsmath,amssymb,amsthm,mathtools,booktabs,array}
\usepackage[numbers,sort&compress]{natbib}
\usepackage[hidelinks]{hyperref}
\usepackage{microtype}
\hypersetup{pdftitle={Marcinkiewicz--Zygmund constants of Clenshaw--Curtis and Padua rules beyond half-exactness},
 pdfauthor={Research development draft},pdfsubject={Complete proof draft; Claude audit and historical novelty checks pending}}
\newtheorem{theorem}{Theorem}[section]
\newtheorem{lemma}[theorem]{Lemma}
\newtheorem{proposition}[theorem]{Proposition}
\newtheorem{corollary}[theorem]{Corollary}
\theoremstyle{remark}
\newtheorem{remark}[theorem]{Remark}
\newcommand{\Pn}{\Pi}
\newcommand{\R}{\mathbb R}
\newcommand{\CC}{\mathrm{CC}}
\newcommand{\Pa}{\mathrm P}
\newcommand{\lamin}{\lambda_{\min}}
\newcommand{\lamax}{\lambda_{\max}}
\newcommand{\norm}[1]{\left\lVert#1\right\rVert}
\newcommand{\dd}{\,\mathrm d}
\newcommand{\id}[1]{\textnormal{\texttt{#1}}}
\DeclareMathOperator{\diag}{diag}
\title{Marcinkiewicz--Zygmund constants of Clenshaw--Curtis\texorpdfstring{\\}{ }and Padua rules beyond half-exactness}
\author{Research development draft\thanks{Authorship and submission venue to be decided by AC. This document is not for submission.}}
\date{31 October 2026}
\begin{document}
\maketitle
\begin{center}\small
\textbf{Development status.} Independent Astra reconstruction accepted the analytic
M1--M3 results within its stated scope. Claude audit~01 accepted the original mathematical statements; focused
re-audit of this response: \textbf{PENDING}.
Historical priority: \textbf{screened, not cleared} (documented, non-exhaustive search; some fulltexts pending); see Section~\ref{sec:context}.
\end{center}

\begin{abstract}
For the Lebesgue Clenshaw--Curtis rule on the $N+1$ Chebyshev extrema,
we prove that its squared-polynomial norm ratio is less than two through
degree $N-1$. An absolute-coefficient comparison for a matrix with
nonpositive off-diagonal entries gives the explicit upper bound
$2-4/N+4/N^2$ by specializing its degree-sensitive form.
At degree $N-1$, weighted estimates for cardinal
interpolation and a parity-compatible harmonic combination of second-kind
Chebyshev polynomials give a minimum Gram eigenvalue of order
$(1+\log N)^{-2}$. Thus the Gram condition number has order
$(1+\log N)^2$, while its deviation from the identity tends to one.
A filtered Fej\'er reconstruction also supplies explicit, nonsharp strong
Marcinkiewicz--Zygmund constants for every fixed degree ratio below one.
The upper bound yields a uniform hyperinterpolation error factor
$(1+\sqrt2)\sqrt V$ under the stated reproduction condition.
An exact witness at $N=512$ excludes a uniform defect bound of $3/4$
covering that endpoint case, without contradicting a conditional theorem
of An and Wu. For signed Lebesgue Padua rules, we prove the exact
factor-two Clenshaw--Curtis marginal and its spectral compression consequence;
the two-dimensional upper bounds and endpoint asymptotics remain open.
All proofs are self-contained. Qualitative failure of uniform converse
Marcinkiewicz--Zygmund inequalities at this borderline, and a positive lower
bound for small degree ratios, follow from classical weighted results. In a
documented but non-exhaustive search we found no prior statement of the
two-sided logarithmic order or of the explicit constants.
\end{abstract}

\section{Problem, conventions, and relation to the literature}\label{sec:context}

Let $Q(f)=\sum_j w_j f(x_j)$ approximate a finite positive measure $\mu$,
and let $G=(Q(P_iP_j))_{ij}$ in a real $L^2(\mu)$-orthonormal basis
of a polynomial space. The basis-independent quantity
\begin{equation}\label{eq:eta}
 \eta=\norm{I-G}_2
 =\max\{1-\lamin(G),\lamax(G)-1\}
\end{equation}
is the smallest nonnegative constant in
\[
 |Q(p^2)-\norm p_{L^2(\mu)}^2|\le\eta\norm p_{L^2(\mu)}^2.
\]
In particular, $\eta<1$ means that every Gram eigenvalue lies in $(0,2)$.
If $Q$ is exact through degree $d$, exactness alone gives $G=I$ through
degree $\lfloor d/2\rfloor$. We call this the half-exactness range.

An, Sommariva, and Vianello (ASV) investigate the role of these constants
and numerically examine their behavior beyond half-exactness for several
rules, including Clenshaw--Curtis (CC) on the interval and Padua rules on
the square \citep{LIT-045}. Their inspected introduction and
Sections~3.1--3.2 distinguish degree-dependent weak constants from strong
constants uniform along a sampling regime. The displayed experiments use
exactness degrees $m=1,\ldots,20$ and approximation degrees
$n=0,\ldots,30$. The present questions are the all-$N$ CC upper threshold,
its endpoint lower spectrum, and uniformity at every fixed ratio $n/N<1$.
These questions are narrower than a classification of all interpolatory rules.

Throughout the CC analysis,
\begin{equation}\label{eq:nodes}
 N\ge2,\qquad \theta_j=\frac{j\pi}{N},\quad x_j=\cos\theta_j,
 \quad 0\le j\le N,
 \qquad Q_N(f)=\sum_{j=0}^N w_j f(x_j),
\end{equation}
where $w_j=\int_{-1}^1\ell_j(x)\dd x$ and $\ell_j\in\Pn_N$ is
the cardinal polynomial at $x_j$. The measure is unnormalized Lebesgue
measure $\dd x$. Write $G_{N,n}$ for the Gram matrix on $\Pn_n$ and
$\lambda_\pm(N,n)$ for its extreme eigenvalues. Thus
\begin{equation}\label{eq:variational}
 \lambda_-(N,n)=\inf_{0\ne p\in\Pn_n}
 \frac{Q_N(p^2)}{\int_{-1}^1p^2},\qquad
 \lambda_+(N,n)=\sup_{0\ne p\in\Pn_n}
 \frac{Q_N(p^2)}{\int_{-1}^1p^2}.
\end{equation}
The node index $N$, point count $N+1$, and maximal algebraic degree of
exactness $d_N$ are distinct. In these conventions
\begin{equation}\label{eq:ade}
 d_N=\begin{cases}N,&N\text{ odd},\\N+1,&N\text{ even}.
 \end{cases}
\end{equation}
ASV use an exactness-degree parameter; An--Wu's CC parameter $m$ counts
points \citep{LIT-046}. We will keep these meanings separate.

For orientation, our principal CC conclusions are
\begin{align}
 &0<\lambda_-(N,n)\le\lambda_+(N,n)
       \le 2-\frac4N+\frac4{N^2}<2, &&0\le n<N,\label{eq:main-upper}\\
 &\lambda_-(N,N-1)\asymp (1+\log N)^{-2},\label{eq:main-endpoint}\\
 &\frac{(1-n/N)^2}{C_*}\le\lambda_-(N,n)
       \le\lambda_+(N,n)\le1+\max\{2n/N-1,0\}^2,
       &&0\le n<N,\label{eq:main-ratio}
\end{align}
with $C_*=12(1+\pi)^4(1+2/\pi)$. All constants are deliberately nonsharp.
The upper constant in \eqref{eq:main-upper} is the endpoint corollary
of the degree-sensitive estimate \eqref{eq:main-ratio}; it equals the
exact quotient one when $N=2$, and the upper inequality is strict for $N\ge3$.
At full degree $N$ the upper threshold fails, as an explicit $T_N$
test shows below. Padua conclusions will be restricted to an exact
compression and a one-sided asymptotic implication.

\paragraph{Attribution and coverage limits.}
The folded Chebyshev moments used below are established prior algebra:
Trefethen's Theorem~5.2 and formulas (5.2)--(5.4) give the mode aliasing
\citep{LIT-055}; selected coefficient-decay results of
Xiang--Bornemann provide related quadrature context \citep{LIT-056}.
ASV equation~(7) already uses the upper-MZ factor
$1+\sqrt{1+\eta}$, together with an aliasing term, and their
equation~(8) separates the inverse-Gram least-squares construction
\citep{LIT-045}. The generic upper-Gram hyperinterpolation mechanism
in Section~\ref{sec:hyper} is consequently known context; its proof is
included to identify precisely the CC input and the required reproduction.

\paragraph{Historical sources inspected.}
Writing $L_N$ for Lagrange interpolation at the $N+1$ nodes, positivity of
the weights and uniqueness of interpolation give
\[
 \sup_{0\ne p\in\Pn_n}\frac{\int_{-1}^1p^2}{Q_N(p^2)}=\frac1{\lambda_-(N,n)},
\]
i.e.\ $\lambda_-(N,n)^{-1/2}$ is the norm of $L_N$ from nodal data of
degree $\le n$ with the CC-weighted $\ell^2$ norm to $L^2(\dd x)$.
Moreover, $\tfrac12\Delta_N(x_j)\le w_j\le5\Delta_N(x_j)$ with
$\Delta_N(x)=\sqrt{1-x^2}/N+N^{-2}$ (from Lemma~\ref{lem:weight-comparison} and the endpoint weights $w_0=w_N\in\{N^{-2},(N^2-1)^{-1}\}$
below), so CC weights may be exchanged with Christoffel-type weights at a
cost of absolute constants. In this language our endpoint result concerns
a borderline converse Marcinkiewicz--Zygmund inequality.
We have inspected the full texts of the following sources; the precise
reading coverage is recorded with the source identifiers.
\begin{itemize}
\item Sloan--Smith's CC product-integration papers \citep{LIT-047,LIT-048}
control the sum of absolute product weights, i.e.\ a $C[-1,1]$ functional
norm, and contain no quadratic-form or nodal $\ell^2$ estimate.
\item The weighted mean-convergence results of Erd\H{o}s--Feldheim,
Askey, Nevai, and Smith's Lobatto theorem
\citep{LIT-061,LIT-062,LIT-063,LIT-065} bound interpolation from
\emph{uniform} data into weighted $L^p$. Nevai's quasi-Lagrange theorem
includes the CC nodes. Uniform data form a stronger input norm than
weighted nodal $\ell^2$, so these results are consistent with, and do not
imply, logarithmic growth of the $\ell^2\to L^2$ norm.
\item Lubinsky's converse inequalities on the full interpolation space
\citep{LIT-064} require an exponent in an open range. For CC nodes,
Lebesgue measure and $p=2$, the nodal polynomial
$\sin\theta\sin N\theta$ sits exactly at the excluded endpoint.
\item Mastroianni--Russo \citep{LIT-293} characterize converse inequalities
at Gauss nodes for generalized Jacobi and Ditzian--Totik weights.
König--Nielsen \citep{LIT-294} do so for same-weight Jacobi nodes.
Xu's sufficient conditions \citep{LIT-296} meet the same obstruction.
At the zeros of $U_{N-1}$ (the interior CC nodes) with Lebesgue measure,
the Mastroianni--Russo condition requires $(1-x^2)^{-1/2}\in L^2$, which
fails logarithmically. Thus the \emph{qualitative} unboundedness behind
$\lambda_-(N,N-1)\to0$ is classical; so is the growth
$\int U_m^2\dd x\asymp\log m$ used in its necessity proofs (compare
Proposition~\ref{prop:counterexample}).
\item The algebraic Marcinkiewicz inequality of Mastroianni--Totik
\citep{LIT-051} applies at CC nodes when $n\le N/(\pi M)$ for their
unspecified constant $M$. Through the weight comparison above it yields
$\lambda_-(N,n)\ge c>0$ for $n\le\theta_0N$, with unspecified $\theta_0$
and $c$. Section~\ref{sec:ratio} adds explicit constants for every fixed
$\theta<1$ and an upper constant below two.
\item The forward quadrature-sum theorem of Lubinsky--M\'at\'e--Nevai
\citep{LIT-050} applies at CC nodes for every degree $\le lN$. It gives
$\lambda_+(N,n)\le C$ with an unspecified absolute constant. The threshold
$2$ and the explicit constants of Theorem~\ref{thm:ratio-upper} and its
endpoint corollary are not implied.
\item The later Sloan--Smith papers \citep{LIT-297,LIT-298} relate
product-integration convergence to mean convergence with uniform data.
\item The interpolation book of Mastroianni--Milovanovi\'c
\citep{LIT-049} (chapters 2, 4 and 5), Lubinsky's survey \citep{LIT-052},
Widom's and Widom--Wilf's Hankel asymptotics \citep{LIT-053,LIT-054},
the Chebyshev--Legendre method \citep{LIT-058} and recent
Chebyshev--Lobatto interpolation estimates in the Legendre norm
\citep{LIT-295} concern uniform norms, Sobolev inputs or different matrix
problems.
\end{itemize}
We found no statement of the two-sided order
$\lambda_-(N,N-1)\asymp(1+\log N)^{-2}$, of the parity-compatible witness
in $\Pn_{N-1}$, of an all-$N$ upper constant below two, or of explicit
fixed-ratio constants. The search is documented but not exhaustive.
The published version of \citep{LIT-045} states that the CC behaviour
beyond half-exactness is, to its authors' knowledge, not known
theoretically. Still unread is a 2000 paper of Mastroianni--Russo on
weighted Marcinkiewicz inequalities. No claim of a first proof is made.
The analytic arguments below do not depend on any of these sources or on
any private proof.

\section{Exact CC identities and an upper bound below two}\label{sec:upper}

Put $T_k(\cos\theta)=\cos(k\theta)$ and
\[
 \mu_k=\int_{-1}^1T_k(x)\dd x
 =\begin{cases}0,&k\text{ odd},\\2/(1-k^2),&k\text{ even}.
 \end{cases}
\]
The moment formula follows directly by integrating
$\sin\theta\cos(k\theta)$; $k=0,1$ can be checked separately.

\begin{lemma}[Exact identities; \id{PD-T013}]\label{lem:identities}
Set $\gamma_0=\gamma_N=1/2$ and $\gamma_j=1$ otherwise. Then
\begin{equation}\label{eq:weights}
 w_j=\frac{2\gamma_j}{N}\left[
 1-\sum_{r=1}^{\lfloor(N-1)/2\rfloor}
      \frac{2\cos(2r\theta_j)}{4r^2-1}
 -\boldsymbol1_{2\mid N}\frac{\cos(N\theta_j)}{N^2-1}
 \right]>0.
\end{equation}
The endpoint weights equal $N^{-2}$ for odd $N$ and $(N^2-1)^{-1}$
for even $N$. For $0\le k\le2N$,
\begin{equation}\label{eq:alias}
 Q_N(T_k)=\begin{cases}\mu_k,&k\le N,\\\mu_{2N-k},&N<k\le2N.
 \end{cases}
\end{equation}
For $0\le i,j\le n\le N$, the continuous and discrete Chebyshev
mass matrices satisfy
\begin{equation}\label{eq:mass}
 M_{ij}=\frac{\mu_{|i-j|}+\mu_{i+j}}2,\quad
 R_{ij}=\frac{\mu_{|i-j|}+Q_N(T_{i+j})}2,\quad
 D:=2(M-R),\quad D_{ij}=\mu_{i+j}-Q_N(T_{i+j}).
\end{equation}
For $n<N$, $D$ is entrywise nonnegative, with positive entries precisely
when $i+j>N$ is even.
\end{lemma}
\begin{proof}
Finite geometric sums give
$\sum_{j=0}^N\gamma_j\cos(k\theta_j)\cos(l\theta_j)
=N\delta_{kl}/(2\gamma_k)$ for $0\le k,l\le N$.
Indeed the corresponding weighted cosine sum vanishes at indices
$1,\ldots,2N-1$ and is $N$ at $0,2N$; the product formula gives the claim.
Thus
\[
 \ell_j(\cos\theta)=\frac{2\gamma_j}{N}
       \sum_{k=0}^N\gamma_k\cos(k\theta_j)\cos(k\theta).
\]
Integration proves \eqref{eq:weights}. The telescoping identity
\begin{equation}\label{eq:telescope}
 \sum_{r=1}^q\frac2{4r^2-1}=1-\frac1{2q+1}
\end{equation}
shows that its bracket is at least $1/N$ for odd $N$ and
$N/(N^2-1)$ for even $N$. It also proves the endpoint formulas.
Cardinal interpolation gives exactness on $\Pn_N$, and
$T_k(x_j)=T_{2N-k}(x_j)$ gives \eqref{eq:alias}.
The identity $T_iT_j=(T_{i+j}+T_{|i-j|})/2$ proves \eqref{eq:mass}.
If $s=i+j>N$ is even and $n<N$, then $r=2N-s$ satisfies $2\le r<s$,
so $D_{ij}=2/(r^2-1)-2/(s^2-1)>0$. Odd moments vanish.
Finally symmetry gives exactness on degree $N+1$ when $N$ is even.
For odd $N$ the moment $N+1$ aliases to $N-1$ and differs from it;
for even $N$ the moment $N+2$ aliases to $N-2$ and differs from it.
This also proves the maximality assertion \eqref{eq:ade}.
\end{proof}

\begin{lemma}[Coefficient mass floor; \id{PD-T015}]\label{lem:floor}
For $n=N-1$ and $E=\diag(2,1,\ldots,1)$,
\begin{equation}\label{eq:floor}
 \frac1N E\preceq M\preceq\frac\pi2 E.
\end{equation}
\end{lemma}
\begin{proof}
The off-diagonal entries of $M$ are nonpositive. Let $b_0=1/2$,
$b_i=1$ for $i\ge1$, and let $m_i$ be the least integer at least $N$
having the parity of $i$. Put $k_i=m_i-1\in\{N-1,N\}$. Telescoping yields
\begin{equation}\label{eq:rowfloor}
 (Mb)_0=\frac1{k_0},\qquad
 (Mb)_i=\frac12\left(\frac1{k_i-i}+\frac1{k_i+i}\right)
       =\frac{k_i}{k_i^2-i^2}\quad(i\ge1).
\end{equation}
To verify these finite row sums, note that the absolutely convergent
even moment series obeys
$\sum_{h\in2\mathbb Z}\mu_{|h|}=2-4\sum_{r\ge1}(4r^2-1)^{-1}=0$.
The infinite row weighted by $b$ is half that sum and is zero.
Its omitted tail, beginning at $j=m_i$ in steps of two, is
\[
 -\sum_{j=m_i,m_i+2,\ldots}
 \left[\frac1{(j-i)^2-1}+\frac1{(j+i)^2-1}\right].
\]
Using $1/(u^2-1)=[1/(u-1)-1/(u+1)]/2$ gives
\eqref{eq:rowfloor}, including $i=0$. Hence
$(Mb)_0/b_0\ge2/N$ and $(Mb)_i/b_i\ge1/N$ for $i\ge1$.
For a real vector $a$, direct expansion gives
\[
 a^TMa=\sum_i\frac{(Mb)_i}{b_i}a_i^2
 +\sum_{i<j}(-M_{ij})b_i b_j
             \left(\frac{a_i}{b_i}-\frac{a_j}{b_j}\right)^2.
\]
This proves the lower bound. If $q(\theta)=\sum_i a_i\cos(i\theta)$,
then $a^TMa=\int_0^\pi q^2\sin\theta\dd\theta
\le\int_0^\pi q^2\dd\theta=(\pi/2)a^TEa$, proving the upper bound.
\end{proof}

\begin{theorem}[Strict threshold and endpoint refinement; \id{PD-T014}, \id{PD-T022}]\label{thm:upper}
For $0\ne p\in\Pn_{N-1}$,
\[
 0<Q_N(p^2)\le\left(2-\frac4N+\frac4{N^2}\right)\int_{-1}^1p^2
 <2\int_{-1}^1p^2.
\]
The same bound holds on every $\Pn_n$ with $n<N$.
For $N=2$ every quotient equals one; for $N\ge3$ the upper inequality is strict.
\end{theorem}
\begin{proof}
Set $A=2M-R=M+D/2$. We show that its off-diagonal entries are
nonpositive. Opposite parities give zero. For distinct equal-parity indices,
write $d=|i-j|\ge2$ and $s=i+j$. If $s\le N$, $A_{ij}=M_{ij}<0$.
Otherwise let $r=2N-s$. Since $i,j\le N-1$,
$r-d=2(N-\max(i,j))\ge2$, and
\begin{equation}\label{eq:zentry}
 A_{ij}=-\frac1{d^2-1}-\frac2{s^2-1}+\frac1{r^2-1}<0.
\end{equation}
Write $p=\sum_i a_iT_i$ and $v_i=|a_i|$. The off-diagonal signs imply
$a^TAa\ge v^TAv$: diagonal terms agree and $a_i a_j\le v_i v_j$.
Since $D$ is entrywise nonnegative,
\[
 2\int p^2-Q_N(p^2)\ge v^TMv+\tfrac12v^TDv\ge v^TMv.
\]
This is positive for $v\ne0$, proving the qualitative threshold of
\id{PD-T014} without the coefficient floor. Positivity of $Q_N(p^2)$ follows from
positive weights and the impossibility that a nonzero polynomial of
degree at most $N$ vanish at all $N+1$ nodes. Restriction proves the rest.

The explicit endpoint refinement is the $n=N-1$ corollary of the independent
degree-sensitive sign comparison in Theorem~\ref{thm:ratio-upper}
(\id{PD-T022}): $\rho=(1-2/N)^2$ and
$1+\rho=2-4/N+4/N^2$. At $N=2$, $\rho=0$ and exactness applies;
at $N\ge3$ the strict case of that theorem applies.
\end{proof}

\begin{remark}[Retained alternative quantitative proof]
The independent mass inequality of Lemma~\ref{lem:floor} (\id{PD-T015})
also gives
$v^TMv\ge N^{-1}v^TEv=N^{-1}a^TEa\ge2(\pi N)^{-1}a^TMa$
in the preceding sign argument, proving the historical bound
\[
 Q_N(p^2)\le\left(2-\frac2{\pi N}\right)\int_{-1}^1p^2.
\]
That proof remains valid, but its constant is quantitatively superseded
by Theorem~\ref{thm:upper}, since $4-4/N\ge2>2/\pi$ for every $N\ge2$.
The mass inequality retains its independent statement and proof.
\end{remark}

\begin{remark}[The sign argument and its boundary]
Entrywise nonnegativity of $D$ does not mean $D\succeq0$.
For $N\ge5$, take $i=\lfloor N/2\rfloor$ and $j=i+2\le N-1$.
Then $D_{ii}=0$, $D_{ij}>0$, and the corresponding principal determinant
is $-D_{ij}^2<0$ (\id{PD-N002}). The proof uses the signs of $A$, not
positive semidefiniteness of the defect. Nor can one extend the theorem
to $\Pn_N$: since $T_N(x_j)^2=1$,
\begin{equation}\label{eq:full}
 \frac{Q_N(T_N^2)}{\int_{-1}^1T_N^2}
 =\frac{2(4N^2-1)}{4N^2-2}>2
 \qquad(\id{PD-T016}).
\end{equation}
Here $\int T_N^2=1-1/(4N^2-1)$ by the product identity.
\end{remark}

\section{The endpoint minimum eigenvalue}\label{sec:endpoint}

Write $H_s=\sum_{r=1}^s r^{-1}$ and
\begin{equation}\label{eq:KN}
 h_N=2H_{N+1}-1,\qquad
 K_N=4\pi^2(\pi+1)^2h_N
       \bigl(h_N+2H_{\lfloor N/2\rfloor+1}\bigr).
\end{equation}

\begin{lemma}[Weight comparison]\label{lem:weight-comparison}
For $1\le j\le N-1$,
\begin{equation}\label{eq:weightcomp}
 \frac{\sin\theta_j}{N}\le w_j\le\frac{5\sin\theta_j}{N}.
\end{equation}
If $m_j=1+\min(j,N-j)$, then $w_j\ge m_j/N^2$ for all $0\le j\le N$.
\end{lemma}
\begin{proof}
The absolutely convergent cosine series
\[
 \frac\pi2\sin\theta=1-2\sum_{r\ge1}
                  \frac{\cos(2r\theta)}{4r^2-1}\quad(0\le\theta\le\pi)
\]
follows by integrating against cosine modes; absolute convergence and
uniqueness of the cosine expansion identify the continuous function.
For example, uniqueness follows from density of polynomials in
$\cos\theta$ among continuous functions on $[0,\pi]$.
Let $S_N(\theta)$ denote the bracket in \eqref{eq:weights}.
Telescoping its tail gives
\[
 |S_N(\theta)-(\pi/2)\sin\theta|
 \le\begin{cases}1/N,&N\text{ odd},\\N/(N^2-1),&N\text{ even}.
 \end{cases}
\]
At interior nodes $\sin\theta_j\ge2/N$. The error is therefore at most
$(2/3)\sin\theta_j$ for all $N\ge2$, including the even worst case $N=2$.
Since $w_j=2S_N(\theta_j)/N$ and
$\pi/2-2/3\ge1/2$, $\pi+4/3<5$, this proves
\eqref{eq:weightcomp}. Moreover
$\sin\theta_j\ge2\min(j,N-j)/N$, so
$w_j\ge2\min(j,N-j)/N^2\ge m_j/N^2$ in the interior.
The endpoint formulas prove the remaining cases.
\end{proof}

\begin{theorem}[Weighted cardinal lower estimate; \id{PD-T017}]\label{thm:cardinal}
For every $p\in\Pn_N$,
\begin{equation}\label{eq:cardinal-result}
 \int_{-1}^1p^2\le K_N Q_N(p^2).
\end{equation}
Consequently $\lambda_-(N,N)\ge K_N^{-1}$, and restriction gives
$\lambda_-(N,N-1)\ge\lambda_-(N,N)\ge K_N^{-1}
\ge c(1+\log N)^{-2}$ for an absolute $c>0$.
\end{theorem}
\begin{proof}
The nodal polynomial $(1-x^2)T_N'(x)$ yields
\[
 \ell_k(\cos\theta)=
 \frac{(-1)^{k+1}\sin\theta\sin(N\theta)}
      {c_kN(\cos\theta-\cos\theta_k)},\qquad
 c_0=c_N=2,\quad c_k=1\ (0<k<N),
\]
with continuous nodal extension. Factor the cosine difference and use
$\sin\theta\le2\sin((\theta+\theta_k)/2)$ and
$|\sin((\theta-\theta_k)/2)|\ge|\theta-\theta_k|/\pi$.
Together with $|\sin(N\theta)|\le\min\{1,N|\theta-\theta_k|\}$,
this gives, for $t=N\theta/\pi$,
\begin{equation}\label{eq:cardinal-decay}
 |\ell_k(\cos\theta)|\le\min\{\pi,|t-k|^{-1}\}
       \le\frac{\pi+1}{1+|t-k|}.
\end{equation}
The first bound includes the continuous value at $t=k$.
Partition $[0,\pi]$ into cells $E_j$ defined by $|t-j|\le1/2$,
truncated at the endpoints. On $E_j$,
\begin{equation}\label{eq:cell}
 |\ell_k(\cos\theta)|\le\frac{D_0}{1+|j-k|},\quad D_0=2(\pi+1),
 \qquad \int_{E_j}\sin\theta\dd\theta\le\frac{\pi^2m_j}{N^2}.
\end{equation}
For the integral, use cell length at most $\pi/N$ and
$\sin\theta\le\pi m_j/N$.

Put $A_{jk}=(1+|j-k|)^{-1}$, $B_{jk}=\sqrt{m_j/m_k}\,A_{jk}$,
and $u_j=m_j^{-1/2}$. Because $m_j\le m_k+|j-k|$,
\begin{align*}
 \sum_k B_{jk}\frac{u_k}{u_j}
 &=m_j\sum_k\frac{A_{jk}}{m_k}
 \le\sum_k A_{jk}+\sum_k\frac1{m_k}
 \le h_N+2H_{\lfloor N/2\rfloor+1},\\
 \sum_j B_{jk}\frac{u_j}{u_k}&=\sum_j A_{jk}\le h_N.
\end{align*}
The harmonic bounds deliberately overcount distances near the endpoints.
The weighted Schur inequality therefore gives
\begin{equation}\label{eq:schur-discrete}
 \norm B_2^2\le h_N\bigl(h_N+2H_{\lfloor N/2\rfloor+1}\bigr).
\end{equation}
For completeness, if the weighted row and column bounds are $R,S$,
then Cauchy--Schwarz gives
$|\sum_kB_{jk}z_k|^2\le(\sum_kB_{jk}u_k)
\sum_kB_{jk}|z_k|^2/u_k$; summing in $j$ gives $RS\sum_k|z_k|^2$.
This proves the stated version without an interpolation theorem.

Set $v_k=p(x_k)$. Cardinal interpolation, \eqref{eq:cell}, and
\eqref{eq:schur-discrete} give
\begin{align*}
 \int p^2
 &\le\frac{\pi^2D_0^2}{N^2}
       \sum_jm_j\left(\sum_kA_{jk}|v_k|\right)^2\\
 &\le K_N\sum_k\frac{m_k}{N^2}|v_k|^2
 \le K_NQ_N(p^2).
\end{align*}
The last inequality is Lemma~\ref{lem:weight-comparison}.
Finally $H_s\le1+\log s$ gives the asserted order.
\end{proof}

The matching upper estimate needs a degree-$N-1$ test polynomial.
Strictly alternating nodal values with positive amplitudes instead represent
a full degree-$N$ interpolant. We enforce both degree and endpoint cancellation
explicitly.

\begin{theorem}[Parity-compatible harmonic witness; \id{PD-T018}]\label{thm:harmonic}
For every $N\ge2$,
\begin{equation}\label{eq:harmonic-result}
 \lambda_-(N,N-1)\le\frac{128000}{(\log N)^2}.
\end{equation}
Together with Theorem~\ref{thm:cardinal}, this proves
\eqref{eq:main-endpoint}.
\end{theorem}
\begin{proof}
Let $U_r(\cos\theta)=\sin((r+1)\theta)/\sin\theta$.
For $N\ge3$ put $M_0=\lfloor(N-1)/2\rfloor$ and
\begin{equation}\label{eq:witness}
 p_N=\sum_{m=1}^{M_0}\frac{
 U_{N-1-2m}-(1-2m/N)U_{N-1}}m.
\end{equation}
Every term has degree at most $N-1$ and parity $N-1$.
The endpoint values $U_r(\pm1)=(\pm1)^r(r+1)$ show that each
bracket vanishes at both endpoints. Define
\[
 B(\theta)=\sum_{m=1}^{M_0}\frac{\sin(2m\theta)}m,\quad
 A(\theta)=\sum_{m=1}^{M_0}\frac{\cos(2m\theta)}m,\quad
 C(\theta)=A(\theta)-H_{M_0}+\frac{2M_0}{N}.
\]
The sine subtraction formula gives the exact identity
\begin{equation}\label{eq:witness-identity}
 p_N(\cos\theta)=
 \frac{\sin(N\theta)C(\theta)-\cos(N\theta)B(\theta)}{\sin\theta}.
\end{equation}

We record elementary bounds on these finite harmonic sums.
For $0<\theta\le\pi/2$ let $K=\lfloor1/\sin\theta\rfloor$.
A geometric sum bounds
$|\sum_{m=a}^b e^{2im\theta}|\le1/\sin\theta$.
Abel summation bounds a harmonic tail from $K+1$ to any $M_0>K$
by $1/((K+1)\sin\theta)\le1$.
The initial segment satisfies
$|\sum_{m\le K}\sin(2m\theta)/m|\le2K\theta\le\pi$
and $|\sum_{m\le K}\cos(2m\theta)/m|\le H_K
\le1+\log(1/\sin\theta)$; truncate it if $M_0\le K$.
Reflection treats the other half of the interval. Thus
\begin{equation}\label{eq:harmonic-bounds}
 |B(\theta)|\le\pi+1<5,
 \quad A(\theta)\le2+\log\frac{\pi}{2\theta}
 \quad(0<\theta\le\pi/2).
\end{equation}
Since $H_{M_0}\ge\log(M_0+1)$, $M_0+1\ge N/2$, and $2M_0/N\le1$,
\begin{equation}\label{eq:Cbound}
 C(\theta)\le3+\log\pi-\log(N\theta)
              \le5-\log(N\theta).
\end{equation}

At interior nodes \eqref{eq:witness-identity} becomes
$p_N(x_k)=-(-1)^kB(\theta_k)/\sin\theta_k$.
The endpoints are zero. Lemma~\ref{lem:weight-comparison} and
$\sin\theta_k\ge2\min(k,N-k)/N$ imply
\begin{equation}\label{eq:witness-discrete}
 Q_N(p_N^2)\le\frac{125}{N}\sum_{k=1}^{N-1}\frac1{\sin\theta_k}
 \le125H_{\lfloor N/2\rfloor}\le125(1+\log N).
\end{equation}

Write $L=\log N$ and first suppose $L\ge128$.
On $E=[N^{-1/2},1/4]$, \eqref{eq:Cbound} gives
$C\le5-L/2\le-L/4$. The elementary inequality
$(a-b)^2\ge a^2/2-b^2$ applied in \eqref{eq:witness-identity} yields
\[
 \int_{-1}^1p_N^2\ge\frac{L^2}{32}I-25J,\quad
 I=\int_E\frac{\sin^2(N\theta)}{\sin\theta}\dd\theta,\quad
 J=\int_E\frac1{\sin\theta}\dd\theta.
\]
Using $\sin\theta\le\theta$,
\[
 I\ge\tfrac12\log(\sqrt N/4)
       -\tfrac12\left|\int_E\frac{\cos(2N\theta)}\theta\dd\theta\right|.
\]
One integration by parts bounds the last absolute integral by $N^{-1/2}$:
the endpoint terms and the integral of $\theta^{-2}$ sum to
$1/(N\cdot N^{-1/2})$. Hence
$I\ge L/4-\log2-1/(2\sqrt N)\ge L/8$.
Since $\sin\theta\ge\theta/2$ on $E$, $J\le2\log(\sqrt N/4)\le L$.
It follows that
\begin{equation}\label{eq:witness-continuous}
 \int p_N^2\ge\frac{L^3}{256}-25L\ge\frac{L^3}{512},
\end{equation}
where $L^2\ge128^2>12800$ suffices. This also proves $p_N\ne0$.
Equations~\eqref{eq:witness-discrete}--\eqref{eq:witness-continuous}
give a quotient at most $250\cdot512/L^2$.
For $0<L<128$, the constant polynomial gives
$\lambda_-(N,N-1)\le1<128000/L^2$.
This covers every remaining $N$, without an unperformed finite check.
\end{proof}

The witness estimate starts at $\log N\ge128$, but
$128000/(\log N)^2<1$ exactly when
$\log N>\sqrt{128000}\approx357.7709$.
Thus this explicit constant is vacuous throughout all computationally feasible
full-spectrum calculations considered here. The finite values in
Table~\ref{tab:cc} do not exhibit or determine an asymptotic leading constant.

\begin{corollary}[Endpoint defect and condition number]\label{cor:endpoint}
For $G_N=G_{N,N-1}$,
\[
 \norm{I-G_N}_2\longrightarrow1,\qquad
 \kappa_2(G_N)\asymp(1+\log N)^2.
\]
\end{corollary}
\begin{proof}
The two endpoint theorems give $\lambda_-(N,N-1)\to0$ and its
two-sided logarithmic order. The constant polynomial has norm ratio one,
so $\lambda_-\le1\le\lambda_+$; Theorem~\ref{thm:upper} gives
$\lambda_+<2$. Therefore
$1-\lambda_-\le\norm{I-G_N}_2<1$, which proves the limit, and
$\lambda_-^{-1}\le\kappa_2(G_N)<2\lambda_-^{-1}$ proves the order.
\end{proof}

\begin{remark}[A rate that has not been proved]\label{rem:rate}
The exact identity
\[
 1-\norm{I-G_N}_2=\min\{\lambda_-(N,N-1),\,2-\lambda_+(N,N-1)\}
\]
limits what follows from the endpoint lower spectrum. The guaranteed gap in
our upper bounds is only of order $N^{-1}$. The arguments do \emph{not} prove
$1-\norm{I-G_N}_2\asymp(\log N)^{-2}$ or eventual dominance of
$1-\lambda_-$ in \eqref{eq:eta}. A sufficiently strong uniform
endpoint upper gap would be additional information.
\end{remark}

\section{Explicit bounds at each fixed degree ratio}\label{sec:ratio}

For sufficiently small ratios $n/N$, a positive lower bound with
unspecified constants already follows from the classical algebraic
Marcinkiewicz inequality of Mastroianni--Totik \citep{LIT-051}
(Section~\ref{sec:context}). The point of this section is explicit
constants for every ratio below one, with an upper constant below two.

\begin{theorem}[Degree-sensitive upper estimate; \id{PD-T022}]\label{thm:ratio-upper}
For $0\le n<N$, put $\rho=\max\{2n/N-1,0\}^2$. Then
\[
 Q_N(p^2)\le(1+\rho)\int_{-1}^1p^2\qquad(p\in\Pn_n).
\]
If $2n\le N$, every quotient equals one. If $2n>N$, the displayed
inequality is strict for every nonzero $p$.
\end{theorem}
\begin{proof}
Exactness settles $2n\le N$. Otherwise set $t=2n/N-1>0$,
$\rho=t^2$, and $A_\rho=(1+\rho)M-R=\rho M+D/2$ on indices $0,\ldots,n$.
Opposite-parity off-diagonal entries vanish; equal-parity unaliased entries
are negative. For an aliased equal-parity pair let
$s=i+j>N$, $d=|i-j|\ge2$, $r=2N-s$, $\delta=N-n>0$.
Then $r<N$, $r-d\ge2\delta$, and
\[
 \frac dr\le1-\frac{2\delta}{N}=t,\qquad
 \frac{d^2-1}{r^2-1}\le\frac{d^2}{r^2}\le\rho.
\]
Consequently
\[
 (A_\rho)_{ij}=-\frac\rho{d^2-1}
                -\frac{1+\rho}{s^2-1}+\frac1{r^2-1}<0.
\]
For this matrix with nonpositive off-diagonal entries, replacement of
coefficients $a$ by $v=|a|$ gives
$a^TA_\rho a\ge\rho v^TMv+v^TDv/2>0$ for $a\ne0$.
This proves the result.
\end{proof}

\begin{corollary}[Endpoint upper constant; \id{PD-T022}]\label{cor:ratio-endpoint}
For every $N\ge2$,
\[
 \lambda_+(N,N-1)\le1+(1-2/N)^2=2-\frac4N+\frac4{N^2}.
\]
Equality holds at $N=2$; the displayed upper inequality is strict for $N\ge3$.
\end{corollary}
\begin{proof}
Set $n=N-1$ in Theorem~\ref{thm:ratio-upper}. Then
$\max\{2n/N-1,0\}=1-2/N$. The case $N=2$ has $2n=N$ and $\rho=0$;
every $N\ge3$ has $2n>N$ and $\rho>0$. This proves the endpoint
refinement used in Theorem~\ref{thm:upper} and \eqref{eq:main-upper}.
\end{proof}

\begin{theorem}[Filtered lower estimate; \id{PD-T023}]\label{thm:ratio-lower}
Put $C_0=2(1+\pi)^2$ and
$C_*=3C_0^2(1+2/\pi)=12(1+\pi)^4(1+2/\pi)$.
For $N\ge2$, $0\le n<N$, and $p\in\Pn_n$,
\begin{equation}\label{eq:fixedlower}
 \int_{-1}^1p^2\le\frac{C_*}{(1-n/N)^2}Q_N(p^2).
\end{equation}
\end{theorem}
\begin{proof}
On the circle $\R/(2\pi\mathbb Z)$ let $d(t,s)\in[0,\pi]$ denote
shorter angular distance, $t_j=j\pi/N$ for $0\le j<2N$,
$q(t)=p(\cos t)$, $\nu(t)=|\sin t|+1/N$, and $\epsilon=1-n/N$.
Integrals on this circle use ordinary $\dd t$; discrete sums use counting
measure unless a factor is written.

For $m\ge1$, the finite Fej\'er identity is
\[
 F_m(t)=\sum_{|k|<m}(1-|k|/m)e^{ikt}
       =\frac1m\left|\sum_{h=0}^{m-1}e^{iht}\right|^2.
\]
Counting pairs with difference $k$ proves the equality. Define
\begin{equation}\label{eq:filter}
 K(t)=\frac{NF_N(t)-nF_n(t)}{N-n},
\end{equation}
where $nF_n=0$ when $n=0$. Its coefficient at frequency $k$ is
$[\max(N-|k|,0)-\max(n-|k|,0)]/(N-n)$.
Thus $K$ has degree $N-1$ and coefficient one for every $|k|\le n$,
including the boundary $|k|=n$. Finite Fourier expansion gives
\begin{equation}\label{eq:reconstruction}
 q(t)=\frac1{2N}\sum_{j=0}^{2N-1}K(t-t_j)q(t_j).
\end{equation}
Indeed the grid average of $e^{i(l-k)t_j}$ vanishes unless $l-k$ is
a multiple of $2N$. Since $|l-k|\le n+N-1<2N$, only $l=k$ remains;
its multiplier is one. No continuous sampling theorem is assumed.

Let $d=d(t,0)$. The geometric-sum formula gives
$mF_m(t)\le m^2$ and $mF_m(t)\le\pi^2/d^2$ because
$\sin(d/2)\ge d/\pi$. The triangle inequality in \eqref{eq:filter}
therefore gives
\begin{equation}\label{eq:filter-decay}
 |K(t)|\le\frac{2N}{\epsilon}\min\{1,\pi^2/(Nd)^2\}
       \le\frac{C_0N}{\epsilon(1+Nd)^2}.
\end{equation}
The last scalar bound follows by splitting $Nd\le\pi$ and $Nd\ge\pi$;
the value at zero follows by continuity.
The function $\nu$ is one-Lipschitz for circle distance and at least $1/N$,
so $\nu(t)/\nu(t_j)\le1+Nd(t,t_j)$. Set
\[
 z_j=\sqrt{\nu(t_j)}q(t_j),\qquad
 H(t,j)=\frac1{2N}K(t-t_j)\sqrt{\frac{\nu(t)}{\nu(t_j)}}.
\]
Then $\sqrt{\nu(t)}q(t)=\sum_jH(t,j)z_j$ and
\begin{equation}\label{eq:Hdecay}
 |H(t,j)|\le\frac{C_0}{2\epsilon}(1+Nd(t,t_j))^{-3/2}.
\end{equation}
Each half-open shell $[k\pi,(k+1)\pi)$ of the scaled distance $Nd$
contains at most two nodes, one in each circle direction. This is because
each corresponding half-open arc has the grid spacing $\pi/N$ as its length.
Hence, using a decreasing-sum integral comparison,
\begin{align}
 \sup_t\sum_j|H(t,j)|
 &\le\frac{C_0}{\epsilon}\sum_{k\ge0}(1+k\pi)^{-3/2}
 \le\frac{C_0}{\epsilon}(1+2/\pi)=:R,\label{eq:rowH}\\
 \sup_j\int_0^{2\pi}|H(t,j)|\dd t
 &\le\frac{C_0}{2\epsilon}\frac2N
                 \int_0^\infty(1+u)^{-3/2}\dd u
 =\frac{2C_0}{\epsilon N}=:S.\label{eq:colH}
\end{align}
At an antipode the node count can only decrease. Cauchy--Schwarz gives
$|\sum_jH(t,j)z_j|^2\le(\sum_j|H(t,j)|)
\sum_j|H(t,j)||z_j|^2$.
Integrating and using \eqref{eq:rowH}--\eqref{eq:colH} proves
\begin{equation}\label{eq:circlebound}
 \int_0^{2\pi}\nu(t)|q(t)|^2\dd t
 \le\frac{2C_0^2(1+2/\pi)}{\epsilon^2N}
          \sum_{j=0}^{2N-1}\nu(t_j)|q(t_j)|^2.
\end{equation}

For each interior half-circle node, $\sin t_j\ge2/N$ and
$\nu(t_j)\le(3/2)\sin t_j$. Lemma~\ref{lem:weight-comparison} and
the endpoint formulas imply at every half-circle node
$w_j\ge2\nu(t_j)/(3N)$. Reflection then gives
\[
 \sum_{j=0}^{2N-1}\nu(t_j)|q(t_j)|^2
 \le2\sum_{j=0}^N\nu(t_j)|p(x_j)|^2\le3NQ_N(p^2).
\]
The first inequality overcounts the two endpoints, which is harmless.
Finally $2\int_{-1}^1p^2=\int_0^{2\pi}|\sin t|q(t)^2\dd t
\le\int_0^{2\pi}\nu(t)q(t)^2\dd t$.
Substitution in \eqref{eq:circlebound} proves \eqref{eq:fixedlower}.
\end{proof}

\begin{corollary}[Explicit strong constants; \id{PD-T024}]\label{cor:strong}
For $0\le\theta<1$, set
\[
 A_\theta=\frac{(1-\theta)^2}{C_*},\qquad
 \rho_\theta=\max\{2\theta-1,0\}^2,\qquad
 \eta_\theta=\max\{1-A_\theta,\rho_\theta\}<1.
\]
For every integer $0\le n\le\theta N$, the spectrum of $G_{N,n}$ lies
in $[A_\theta,1+\rho_\theta]$; consequently
\[
 (1-\eta_\theta)\int p^2\le Q_N(p^2)
                    \le(1+\eta_\theta)\int p^2\qquad(p\in\Pn_n).
\]
\end{corollary}
\begin{proof}
The two preceding theorems give the spectral interval by monotonicity in
$n/N$. Since $C_*>1$, $0<A_\theta<1$ and $0\le\rho_\theta<1$.
The definition of $\eta_\theta$ gives both inequalities.
\end{proof}

The filtered estimate degenerates like $N^{-2}$ at $n=N-1$ and does not
replace the sharper endpoint theorem. Conversely the endpoint cardinal
argument does not supply a constant independent of $N$ at a fixed ratio.
Both proofs are included because the two regimes need different estimates.

\section{The three-quarter assumption and hyperinterpolation}\label{sec:hyper}

An--Wu's inspected arXiv version assumes a relative squared-norm MZ error
$\eta<1$ in its main theorem. Its Section~4 assumes $\eta=3/4$ for the
examples, and Corollary~4.1 retains the theorem's conditions for an
$m$-point CC rule with $n+2\le m\le2n+1$, obtaining an error bound
$3\sqrt V E_{m-1-n}(f)$ \citep{LIT-046}.
The following finite calculation specifies an excluded case for that
assumption; it does not refute the conditional implication.

\begin{proposition}[An exact endpoint witness; \id{PD-T019}, \id{PD-N004}]\label{prop:counterexample}
For $S_N=\sum_{j=1}^N(2j-1)^{-1}$,
\begin{equation}\label{eq:Uexact}
 \int_{-1}^1U_{N-1}^2=2S_N,\qquad
 Q_N(U_{N-1}^2)=
 \begin{cases}2,&N\text{ odd},\\2N^2/(N^2-1),&N\text{ even}.
 \end{cases}
\end{equation}
In particular $\norm{I-G_{512,511}}_2>3/4$. This case has $513$ CC
points, so An--Wu's point-count notation is $m=513$, $n=511$.
\end{proposition}
\begin{proof}
The finite trigonometric product identity
$U_r^2=\sum_{s=0}^rU_{2s}$ follows by summing the geometric sine series.
Also $\int_{-1}^1U_{2s}=\int_0^\pi\sin((2s+1)\theta)\dd\theta
=2/(2s+1)$, giving the continuous formula.
All interior nodal values of $U_{N-1}$ vanish; its endpoint values have
absolute value $N$. The endpoint weights prove the discrete formula.

For a short exact certificate at $N=512$, direct rational addition gives
$S_8>201/100$. Monotonicity of $1/(2t-1)$ gives
\[
 S_{512}\ge S_8+\int_9^{513}\frac{\dd t}{2t-1}
          =S_8+\tfrac12\log(1025/17).
\]
The factorial series for $e$ gives
$e\le1+1+1/2+1/6+(1/24)/(1-1/5)=87/32<11/4$.
The rational inequality $(11/4)^4<1025/17$ implies
$\log(1025/17)>4$, so $S_{512}>401/100$.
Since $4\cdot512^2/(512^2-1)<401/100$, \eqref{eq:Uexact} yields
\[
 \frac{Q_{512}(U_{511}^2)}{\int U_{511}^2}
 =\frac{512^2}{(512^2-1)S_{512}}<\frac14.
\]
Thus $\lambda_-(512,511)<1/4$ and the defect exceeds $3/4$.
\end{proof}

The growth $\int_{-1}^1U_m^2\asymp\log m$ behind this witness is
classical; it drives the necessity part of the converse
Marcinkiewicz--Zygmund characterizations \citep{LIT-293,LIT-294}. The
new content here is only the exact CC evaluation.
This simple witness gives only an $O(1/\log N)$ quotient asymptotically;
the harmonic combination in Theorem~\ref{thm:harmonic} is needed for
the squared logarithm. The finite certificate uses no floating spectrum.
The published An--Wu Corollary~4.1 coincides with the inspected arXiv
formulation: it assumes $\eta=3/4$ for $m$-point CC with $n+2\le m\le2n+1$
\citep{LIT-046}, and $m=513$, $n=511$ lies in that range. No correction
was found in a documented search; this is not proof of absence.

\begin{proposition}[Upper-Gram hyperinterpolation estimate; \id{PD-T025}]\label{prop:hyper}
Let $\mu$ be a finite positive measure of mass $V>0$ on a compact domain.
Let $Q$ have positive weights of total mass $V$, and let
$P_1,\ldots,P_d$ be an $L^2(\mu)$-orthonormal basis of $\Pn_n$.
Assume $Q(p^2)\le B\norm p_2^2$ on $\Pn_n$, and define
\[
 \mathcal L_nf=\sum_{i=1}^d Q(fP_i)P_i.
\]
Then $\norm{\mathcal L_nf}_2\le\sqrt{VB}\norm f_\infty$.
If $0\le k\le n$ and $Q$ is exact on all products of a member of
$\Pn_k$ with a member of $\Pn_n$, then
\begin{equation}\label{eq:hypererror}
 \norm{f-\mathcal L_nf}_2\le\sqrt V(1+\sqrt B)E_k(f),\qquad
 E_k(f)=\inf_{r\in\Pn_k}\norm{f-r}_\infty.
\end{equation}
\end{proposition}
\begin{proof}
For $c_i=Q(fP_i)$ and any Euclidean unit vector $a$, put
$p_a=\sum_i a_iP_i$, so $\norm{p_a}_2=1$.
Weighted Cauchy--Schwarz gives
$|a\cdot c|=|Q(fp_a)|\le Q(f^2)^{1/2}Q(p_a^2)^{1/2}
\le\sqrt{VB}\norm f_\infty$.
Euclidean duality and orthonormality identify
$\norm{\mathcal L_nf}_2=\norm c_2$, proving the first bound.
For $r\in\Pn_k$, product exactness gives
$Q(rP_i)=\int rP_i\dd\mu$, and since $k\le n$,
$\mathcal L_nr=r$. Hence
\[
 \norm{f-\mathcal L_nf}_2
 \le\norm{f-r}_2+\norm{\mathcal L_n(f-r)}_2
 \le\sqrt V(1+\sqrt B)\norm{f-r}_\infty.
\]
Taking the infimum proves \eqref{eq:hypererror}.
\end{proof}

For CC, $V=2$, $0\le n<N$, and one may take
\begin{equation}\label{eq:Bcc}
 B_{N,n}=1+\rho_{N,n}=1+\max\{2n/N-1,0\}^2<2.
\end{equation}
Indeed $n\le N-1$ gives $1+\rho_{N,n}\le2-4/N+4/N^2<2-2/(\pi N)$;
the minimum of the two earlier bounds always selects $1+\rho_{N,n}$.
Thus for every $0\le k\le\min(n,d_N-n)$,
\begin{equation}\label{eq:CCerror}
 \norm{f-\mathcal L_nf}_{L^2(-1,1)}
 \le\sqrt2(1+\sqrt{B_{N,n}})E_k(f)
 \le(1+\sqrt2)\sqrt2\,E_k(f).
\end{equation}
Using only exactness through $N$ permits $k\le\min(n,N-n)$;
even-$N$ parity can improve that by one when $k\le n$ allows it.
For the $513$-point, degree-$511$ example, $k=1$ already gives the
uniform factor $(1+\sqrt2)\sqrt V$ in place of the assumed
$3\sqrt V$ factor; parity also permits $k=2$ here.
This estimate concerns the defined quadrature hyperinterpolant. An
inverse-Gram least-squares operator is a different construction; growth
of $\lambda_-^{-1/2}$ does not describe the norm of $\mathcal L_n$.
Neither the general mechanism nor the validity of An--Wu's conditional
theorem is challenged by this CC application.

\section{An exact Padua compression}\label{sec:padua}

Let $\mathcal P_n$ denote bivariate polynomials of total degree at most $n$.
For a construction index $N\ge2$, use the Padua support
\begin{equation}\label{eq:padua}
 \mathcal X_N=\left\{\left(\cos\frac{j\pi}{N},
                   \cos\frac{k\pi}{N+1}\right):
    0\le j\le N,\ 0\le k\le N+1,\ j+k\text{ even}\right\}.
\end{equation}
Suppose real weights $W_{jk}$ define a rule $Q_N^\Pa$ exact on
$\mathcal P_N$ for unnormalized Lebesgue measure $\dd x\dd y$ on the
square. This assumption applies to the interpolatory rule obtained by
integrating its interpolation basis or solving its exact moment equations.
The following implication needs neither positivity of the weights nor a
separate existence or uniqueness theorem for that construction.

\begin{theorem}[Marginal and compression; \id{PD-T021}]\label{thm:padua}
At each $x$-coordinate the total weight equals
\begin{equation}\label{eq:marginal}
 \sum_{k:j+k\,\mathrm{even}}W_{jk}=2w_j,\qquad0\le j\le N.
\end{equation}
Let $G^\Pa_{N,n}$ be the real symmetric Gram matrix on $\mathcal P_n$
in a Lebesgue-orthonormal basis. For $0\le n\le N$,
\begin{align}
 \lamin(G^\Pa_{N,n})&\le\lambda_-(N,n),&
 \lamax(G^\Pa_{N,n})&\ge\lambda_+(N,n),\label{eq:padua-extreme}\\
 \norm{I-G^\Pa_{N,n}}_2&\ge\norm{I-G_{N,n}}_2.&&\label{eq:padua-eta}
\end{align}
In particular,
$\liminf_{N\to\infty}\norm{I-G^\Pa_{N,N-1}}_2\ge1$,
and the Padua defect at $N=512,n=511$ exceeds $3/4$.
\end{theorem}
\begin{proof}
Write $a_j=\sum_kW_{jk}$. Exactness on univariate $q(x)\in\Pn_N$ gives
$\sum_j a_jq(x_j)=2\int_{-1}^1q=\sum_j2w_jq(x_j)$.
Taking $q=\ell_j$ proves \eqref{eq:marginal}.
Define $Jq(x,y)=q(x)/\sqrt2$. Then
$\int_{[-1,1]^2}(Jq)^2=\int_{-1}^1q^2$, so $J$ is an isometry from
$\Pn_n$ into $\mathcal P_n$. The exact marginal gives for $p,q\in\Pn_n$
\[
 Q_N^\Pa((Jp)(Jq))=\tfrac12\sum_j a_jp(x_j)q(x_j)=Q_N(pq).
\]
This step uses the marginal rather than cubature exactness for the product,
whose degree can exceed $N$. In orthonormal coordinates,
$J^*G^\Pa_{N,n}J=G_{N,n}$ and $J^*J=I$.
The symmetric variational principle proves \eqref{eq:padua-extreme},
and compression of $I-G^\Pa_{N,n}$ proves \eqref{eq:padua-eta}.
No sign condition on $W_{jk}$ was used.
The endpoint lower eigenvalue tends to zero. Therefore
$\norm{I-G^\Pa_{N,N-1}}_2\ge1-\lambda_-(N,N-1)$ gives the stated
liminf after taking limits on the right.
Proposition~\ref{prop:counterexample} gives the finite consequence.
\end{proof}

The theorem provides no limsup bound, no Padua positive definiteness, and
no two-dimensional upper threshold. In particular it does not prove that
the Padua defect tends to one or that its minimum eigenvalue has the CC
logarithmic order. For signed weights these require a distinct argument
on the complement of the univariate subspace. These targets remain open
(\id{PD-T028}).

\section{Finite observations and unresolved sharpenings}\label{sec:numerics}

Tables~\ref{tab:cc}--\ref{tab:padua} reproduce bounded float64 observations
(\id{PD-T026}), not asymptotic proofs. The producer runs with Python,
NumPy, and SciPy and pins its code and supplied input hashes in the evidence
manifest. An independent Astra agent separately reran the bounded suite.
On the checked CC cases $N=3,4,5,9,10,21,40,81$, each at
$n=\lfloor N/2\rfloor,N-1,N$, the generalized Chebyshev pencil and a
separately assembled Legendre Gram matrix agreed in extreme eigenvalues
within $1.35\cdot10^{-13}$. This validates only the recorded numerical
scope. The analytic theorems above do not use this comparison as proof.

\begin{table}[htbp]
\centering
\caption{Measured CC spectra on $\Pn_{N-1}$, using $N+1$ points.
Values are rounded from the reproduced output.}\label{tab:cc}
\begin{tabular}{rccc}\toprule
$N$ & $\lambda_-$ & $\lambda_+$ & $\eta$\\\midrule
40   & .259849025 & 1.004471121 & .740150975\\
80   & .216153284 & 1.004439735 & .783846716\\
160  & .182490277 & 1.004439101 & .817509723\\
320  & .156015184 & 1.004437801 & .843984816\\
640  & .134833488 & 1.004437222 & .865166512\\
1280 & .117635156 & 1.004437035 & .882364844\\\bottomrule
\end{tabular}
\end{table}

\begin{table}[htbp]
\centering
\caption{Measured signed Padua spectra on total degree $n=m-1$;
$m$ is the Padua construction degree in this table.}\label{tab:padua}
\begin{tabular}{rrcccc}\toprule
$m$ & $n$ & $\lambda_-$ & $\lambda_+$ & $\eta$ & negative weights\\\midrule
10 & 9  & .356191383 & 1.634856615 & .643808617 & 2\\
20 & 19 & .272233005 & 1.723313732 & .727766995 & 2\\
30 & 29 & .236705530 & 1.757929856 & .763294470 & 2\\
40 & 39 & .214817029 & 1.778088624 & .785182971 & 2\\\bottomrule
\end{tabular}
\end{table}

The separately checked Padua marginals at construction degrees $10,20,31$
(covering odd and even degrees)
obeyed $\max_j|a_j-2w_j|<8.73\cdot10^{-16}$.
The two negative weights in each listed Padua experiment underscore the
need to retain a signed-form argument. Finite positive eigenvalues alone
give no all-degree positive-definiteness theorem.

Several sharper questions remain separate from the proved comparisons.
First, neither existence nor value of
$\lim_{N\to\infty}(\log N)^2\lambda_-(N,N-1)$ is proved
(\id{PD-T020}). Second, the apparent upper limit near $1.004437$,
its existence and exact value, and a uniform endpoint bound $B<2$ are
open (\id{PD-T027}). Table~\ref{tab:cc} stops at $N=1280$;
the supplied $N=20480$ output has not been reproduced and supplies no
claim here. Third, the Padua positivity, upper threshold, endpoint order,
and defect convergence remain open as specified after
Theorem~\ref{thm:padua}. These mathematical questions are distinct from
the pending historical source comparison and required Claude audit.

\paragraph{Reproducibility and claim identification.}
The accompanying README gives exact compilation and numerical commands,
artifact hashes, and links to the theory and evidence. The final
\id{PD-T}/\id{PD-N} identifiers in theorem headings refer to the durable
repository claims; bibliography keys use the final \id{LIT} allocation.
These identifiers do not confer publication or audit readiness. All
proofs of the central CC and Padua statements have been included here;
no private manuscript is a mathematical dependency.

\bibliographystyle{plainnat}
\bibliography{refs}
\end{document}
